\subsection{A Lisp Reading Guide}
\label{sec:lisp-reading}

This section teaches the Lisp notation used by the original examples.
Try the small expressions at the ABCL/R2 prompt. Comments beginning with
a semicolon explain the expected result; do not type the prompt itself.
The guide is deliberately selective: it is not a Common Lisp course.

\subsubsection{Expressions, Symbols, and Literal Data}

An expression such as \code{(+ 2 3)} puts the operation first.
For an ordinary function call, argument expressions are evaluated before
the function is called. Numbers and strings are literal values.
An unquoted ordinary symbol in an expression names a binding:
\begin{lstlisting}
(+ 2 (* 3 4))                  ; => 14
(let ((amount 10)) (+ amount 2)) ; => 12
\end{lstlisting}
The inner multiplication finishes before the outer addition.
\code{let} introduces a local binding; its body can refer to \code{amount}.
A \emph{binding} associates a name with a value. Its \emph{lexical scope}
is determined by the surrounding source structure: the body of this
\code{let} is where its local name is available.

A symbol is not a string. \code{amount} is a name, \code{'amount} is that
name as data, and \code{"amount"} is a string of characters.
The quote character is shorthand for \code{quote}:
\begin{lstlisting}
'amount                       ; => AMOUNT
'(a b c)                      ; => (A B C)
(list 'a 'b 'c)               ; => (A B C)
\end{lstlisting}
Quote prevents evaluation of its operand. Thus the symbols inside
\code{'(a b c)} are not looked up as variables.
Unescaped symbol names are normally printed in uppercase; \code{n} and
\code{N} read as the same symbol in the same package.

A \emph{keyword} begins with a colon and evaluates to itself.
It is convenient as a protocol tag: \code{:deposit}, \code{:new},
and \code{:finish} are distinct symbols. No variable definition is
needed to use them.
\begin{lstlisting}
(list :deposit 10)            ; => (:DEPOSIT 10)
(let ((x 7)) [:point x (+ x 1)]) ; => (:POINT 7 8)
\end{lstlisting}
Here square brackets construct a list of evaluated elements.
They also introduce ABCL/R2 operations such as object creation or sending;
the recognized operator determines which interpretation applies.
In a \emph{pattern}, however, a symbol such as \code{x} binds a part
of the incoming value. Construction and matching are different uses
of the same readable list notation.

\subsubsection{Truth, Tests, and Control Flow}

\code{nil} means false and is also the empty list. \code{t} is true.
Every value other than \code{nil}, including zero, counts as true.
\begin{lstlisting}
(if nil :yes :no)             ; => :NO
(if 0 :yes :no)               ; => :YES
(null '())                    ; => T
(not t)                       ; => NIL
\end{lstlisting}
\code{if} is a \emph{special form}: it does not eagerly evaluate both
branches. It evaluates the test and only the selected branch.
\code{progn} evaluates several expressions in sequence, with the final
expression providing its value. \code{when} is a conditional sequence;
\code{cond} chooses the first clause whose test succeeds.
\code{and} stops at the first false operand; \code{or} stops at the
first true operand. These are control operations, not ordinary eager calls.

Use \code{=} to compare numbers. Use \code{eq} for symbol or object
identity, as in testing a terminal keyword or checking a self-transfer.
Use \code{equal} for structural comparison of lists and strings.
Do not assume that \code{eq} compares two independently constructed
lists by their contents.

\subsubsection{Lists, Cons Cells, and Dotted Tails}

A \emph{cons cell} has two parts: its \code{car} and its \code{cdr}.
A proper list chains cons cells and ends in \code{nil}.
\code{first} and \code{rest} are readable names for the head and tail
of a list; \code{car} and \code{cdr} expose the underlying pair.
\begin{lstlisting}
(cons 'a '(b c))              ; => (A B C)
(car '(a b c))                ; => A
(cdr '(a b c))                ; => (B C)
(cons 'd 'e)                  ; => (D . E)
(cdr '(d . e))                ; => E
\end{lstlisting}
The dot means that the final tail is the value after the dot.
\code{(d . e)} is a pair, not the two-element proper list \code{(d e)}.
Its tail is the symbol \code{e}, not the list \code{(e)}.
\code{(a b)} is shorthand for \code{(a . (b . nil))}.

In a bracket pattern, \code{[A . B]} binds the head to \code{A} and
the entire tail to \code{B}. \code{[[A] . B]} additionally requires
the head to be a one-element proper list. These distinctions drive the
original tree walker; do not replace dotted pairs by ordinary lists.

\code{append} concatenates lists, copying all but the final argument.
\code{nconc} performs destructive concatenation.
\code{length} counts elements of a proper sequence; \code{nth} selects
a zero-based element of a list. Later original programs also use arrays:
\code{make-array} creates an array and \code{aref} accesses an element.
Do not pass an arbitrary dotted pair where a proper list is required.

\subsubsection{Bindings and Updating Values}

\code{let} evaluates its initializers before its new bindings are available
to the body. \code{let*} introduces bindings one at a time, so a later
initializer can refer to an earlier binding.
\begin{lstlisting}
(let* ((x 3) (y (+ x 1))) (* x y)) ; => 12
(let ((x 3)) [x := (+ x 1)] x)    ; => 4
\end{lstlisting}
\code{[x := value]} changes an ABCL/R2 binding.
\code{setq} is also available, but this parser permits only one assignment
pair in a \code{setq} form.
\code{setf} is a host macro for updating a \emph{place}, such as the
\code{cdr} of a pair or an array element. Do not confuse changing a binding
with changing an existing data structure.

An \emph{association list} stores bindings as pairs:
\begin{lstlisting}
(let ((bindings '((n . 3) (balance . 20))))
  (cdr (assoc 'balance bindings))) ; => 20
\end{lstlisting}
\code{assoc} finds the pair whose key matches.
\code{pairlis} pairs a list of names with a list of values.
The original environment and script examples use association lists to
describe bindings; those lists are data, whereas an environment is an
object that implements lookup and update messages.
A local set of bindings for one scope is called a \emph{frame}; an
environment can link its frame to an enclosing environment.

\subsubsection{Reading a Loop}

ABCL/R2 implements \code{do} directly. Read its three parts as local
bindings, a termination clause, and a body:
\begin{lstlisting}
(do ((n 0 (+ n 1)))
    ((= n 3) n)
  (format t "~A~%" n))
; Prints 0, 1, 2 on separate lines; returns 3.
\end{lstlisting}
\code{n} starts at zero. Before each iteration the termination test is
checked. If false, the body runs and the step expression updates \code{n}.
If true, the result expressions of the termination clause run.
With several step expressions, their new values are computed before the
loop variables are updated. A variable with no step expression keeps its
value unless the body changes it.

The tree walker uses \code{(do ((tree received-tree)) ((null tree)) ...)}.
There is no step expression: its \code{match} clauses explicitly update
\code{tree}. Its loop result is \code{nil}, but the handler subsequently
sends \code{!:finish}. The loop result is not that reply.

\subsubsection{Functions, Macros, and Reader Notation}

\code{defun} defines a host Lisp function; its parameter list describes
arguments and its body computes a result.
\code{\#'car} is shorthand for \code{(function car)} and denotes a function,
not the symbol \code{'car}. \code{funcall} calls a function value;
\code{mapcar} applies a function to corresponding list elements:
\begin{lstlisting}
(mapcar #'car '((a . 1) (b . 2))) ; => (A B)
\end{lstlisting}
\code{apply} calls a function using a list for its final arguments.
The priority queues use \code{(apply \#'sort item-list predicate keys)}:
when \code{keys} is \code{(:key \#'first)}, this is a call to
\code{sort} with that keyword pair appended to its arguments.
\code{sort} can rearrange its input list destructively; use its returned
list rather than assuming that the old head still names the result.
The host function is called with values, not ABCL forms to execute later.

\begin{lstlisting}
(apply #'+ '(1 2 3))             ; => 6
(let ((items (list :a :b)))
  (push :new items)
  (list (pop items) items))       ; => (:NEW (:A :B))
\end{lstlisting}

\code{push} changes a place to hold a new cons cell; \code{pop} returns
its first element and changes the place to the remaining tail.
\code{pushnew} inserts only when the item is not already present
under its equality test. These macros update bindings or places;
they are not just pure list constructors. \code{caar}, \code{cadr},
and \code{cddr} abbreviate \code{car} of \code{car}, \code{car} of
\code{cdr}, and \code{cdr} of \code{cdr}, respectively.

\code{lambda} describes an anonymous host function. As explained below,
it does not automatically capture ABCL/R2 state.
A \emph{closure} is a function value together with the lexical bindings
it retains. Host Lisp closures retain host Lisp bindings, not arbitrary
bindings in the ABCL/R2 interpreter's environment objects.

A macro transforms source expressions before the transformed program
is translated or executed. Unlike a function, it receives argument forms,
not their computed values. \code{defmacro} defines such a transformation.
Backquote builds a source template; comma inserts a form into it,
and comma-at splices a list of forms into the surrounding list.
The original inheritance macro is a compact example:
\wholelisting[firstline=14,lastline=15]{examples/misc/inherit.abclr2}
\code{\&rest args} collects all argument forms into a list.
Thus \code{(inherit Any Reply Sender bird)} becomes
\code{(reify-args-to-meta :inherit Any Reply Sender bird)}.
The macro inserts the tag and preserves the argument forms; it does not
itself perform inheritance.
You only need to recognize the transformation here;
Section~\ref{sec:reflection} explains what the reflective operation does.

Other parameter markers you will meet are \code{\&optional} for optional
positional arguments and \code{\&key} for keyword arguments.
\code{\&aux} introduces initialized local variables in a host function
or macro's parameter list. The macros in Copy and Report use it with
\code{gensym}, which creates fresh symbols so generated temporary
bindings do not accidentally capture a name in the caller's source.
In \code{(inc "demo" :use-source t)}, the keyword selects an option
and \code{t} supplies its value.
\code{r2sys::create-primary-eval} names an internal symbol in the runtime
package. A package groups symbol names; the double colon exposes an
internal name rather than promising a portable public API.
\code{\#.} in a source file requests evaluation while reading;
it is different from an ordinary expression executed later.
Do not treat a source file containing it as inert quoted data.
\code{\#\textbackslash{}>} is a character literal, distinct from the
one-character string \code{">"}; Copy's display-name helpers inspect
such characters in strings without changing the object's identity.

\code{\#| ... |\#} encloses a block comment; many original files put
sample sessions there, so loading the file does not run those commands.
\code{\#+feature} includes the next form only if that reader feature is
enabled, and \code{\#-feature} does the opposite. The original sources
use \code{\#+COMMENT} for inactive alternatives in the normal build.
This is a reader condition, not a block comment or an extra runtime
branch: do not combine such alternatives with the active implementation.

\code{format} uses a control string: \code{\~{}A} prints a value readably
for a person, \code{\~{}S} prints Lisp notation, and \code{\~{}\%} emits
a newline. The first argument \code{t} selects standard output.
\code{\~{}D} prints a decimal integer and \code{\~{}\&} begins a fresh
line when needed. An open output stream can replace \code{t}; the
monitoring examples use \code{open} to create their log streams and
\code{close} when recording ends.
These explanations are enough to read the printer in the original
future example without assuming familiarity with Common Lisp formatting.

\paragraph{Checkpoint, with answers.}
What are the values of \code{(cdr '(a b))} and \code{(cdr '(a . b))}?
The first is the list \code{(B)}; the second is the symbol \code{B}.
Why does \code{'n} work as an environment lookup key?
Because the environment needs the variable's name as data, not its value.
Why is the inheritance macro not an ordinary function?
Because it constructs a reflective expression from unevaluated source forms.
